\section{Πρόβλημα 1}

\subsection{Ερώτηση (i)}

Μας δίνεται παίγνιο συμφόρησης στο ακόλουθο δίκτυο όπου οι παίκτες είναι
απειροστά μικροί και δημιουργούν συνολική απαίτηση 3 μονάδων. Θα βρούμε το
τίμημα της αναρχίας.

\begin{center}
\begin{tikzpicture}[
  node distance=2.5cm,
  every node/.style={circle, draw, fill=yellow!70!orange, minimum size=0.9cm, font=\bfseries},
  >={Stealth[length=6pt]}
]
  \node (s) {s};
  \node (m) [right=of s] {m};
  \node (t) [right=of m] {t};

  \draw[->] (s) -- node[above, draw=none, fill=none, font=\normalfont] {$1$} (m);
  \draw[->] (m) -- node[above, draw=none, fill=none, font=\normalfont] {$1$} (t);
  \draw[->] (s) to[bend left=50] node[above, draw=none, fill=none, font=\normalfont] {$x$} (t);
  \draw[->] (s) to[bend right=40] node[below, draw=none, fill=none, font=\normalfont] {$x$} (m);

  \draw[->] ($(s)+(-1,0)$) -- node[above, draw=none, fill=none, font=\normalfont] {$3$} (s);
\end{tikzpicture}
\end{center}

Η ακόλουθη ροή αποτελεί σημείο ισορροπίας: $2$ μονάδες στο μονοπάτι
$P_1 = [(s, t)]$ και $1$ μονάδα στο μονοπάτι $P_2 = [(s, m), (m, t)]$ μέσω της
ακμής με κόστος $\ell_{s \to m}(x) = x$. Όλα τα μονοπάτια από το $s$ στο $t$
έχουν κόστος $2$ με αυτό τον τρόπο. Τα σημεία ισορροπίας σε παίγνια συμφόρησης
απειροστά μικρών παικτών έχουν μοναδικά σημεία ισορροπίας, οπότε δεν χρειάζεται
να ψάξουμε για άλλο ``καλύτερο'' σημείο. Το μέσο κόστος της ροής είναι:
$$
    2 \cdot 2 + 1 \cdot (1 + 1) = 6
$$

Για να βρούμε την βέλτιστη ροή, θέτουμε συναρτήσης κόστους
$\ell_{s \to t}'(x) = \ell_{s \to m}'(x) = 2x$ και βρίσκουμε το καινούριο σημείο
ισορροπίας. Το καινούριο σημείο ισορροπίας είναι: $1$ μονάδα στο μονοπάτι
$P_1 = [(s, t)]$, $1.5$ μονάδα στο μονοπάτι $P_2 = [(s, m), (m, t)]$ μέσω της
της ακμής με κόστος $\ell_{s \to m}(x) = 1$ και $0.5$ μονάδα στο μονοπάτι
$P_3 = [(s, m), (m, t)]$ μέσω της ακμής με (καινούριο) κόστος
$\ell_{s \to m}'(x) = 2x$. Η ροή αποτελεί σημείο ισορροπίας αφού όλα τα μονοπάτια
έχουν κόστος $2$. Στον αρχικό γράφο, αυτό θα μας έδινε μέσο κόστος
$$
    1 \cdot 1 + 1.5 \cdot (1 + 1) + 0.5 \cdot (0.5 + 1) = 1 + 3 + 0.75 = 4.75 \,.
$$

Το τίμημα της αναρχίας είναι $ \frac{6}{4.75} $.

\subsection{Ερώτηση (ii)}

Μας δίνεται το ίδιο παίγνιο συμφόρησης στο ίδιο δίκτυο με 3 ατομικούς παίκτες.
Θα βρούμε το τίμημα της αναρχίας.

Χρειάζεται πρώτα να βρούμε το ``χειρότερο'' σημείο ισορροπίας. Παρατηρούμε ότι
κατανέμοντας $1$ παίκτη σε κάθε μονοπάτι θα έχουμε σημείο ισορροπίας με μέσο
κόστος $5$. Ο μόνος τρόπος να πάρουμε μεγαλύτερο κόστος $5$ είναι να μην
χρησιμοποιήσουμε το μονοπάτι σταθερού κόστους. Επίσης, ο μόνος τρόπος να πάρουμε
παραπάνω από $6$ είναι να χρησιμοποιήσουν και οι $3$ παίκτες μία κοινή ακμή από
αυτές με κόστος $\ell(x) = x$, το οποίο όμως δεν είναι σημείο ισορροπίας (ούτε
το $(3, 0, 0)$ ούτε το $(0, 0, 3)$). Αυτό μας δίνει τα χειρότερα σημεία
ισορροπίας $(2, 1, 0), (2, 0, 1), (1, 2, 0)$ και $(1, 0, 2)$ με κόστος $6$.

Για το καλύτερο κοινωνικό κόστος, ήδη γνωρίζουμε ότι υπάρχει η
κατανομή $(1, 1, 1)$ με κόστος $5$. Μη χρησιμοποιώντας την διαδρομή σταθερού
κόστους μπορούμε να επιτύχουμε μόνο μεγαλύτερο κόστος. Βάζοντας περισσότερα
άτομα στην διαδρομή μπορούμε να πετύχουμε το ίδιο ή χειρότερο. Οπότε τελικά το
τίμημα της αναρχίας είναι $\frac{6}{5}$.
